% ====================================================================
% fig_ff1_1_hysgap.tex — [FF1 후보 1] 히스테리시스 gap 닫힌 꼴의 전 거동
%   대상 문서: Ch1 (ch1_preamble 라이브러리: arrows.meta·calc·backgrounds 만 사용)
%   제안 배치: §4.2 식 (eq:dUhys) 문단 뒤(그림 fig:hysloop 다음), §13.5 에서 재참조
%   좌표 생성: eval_ff1_coords.py §P1 (식 eq:dUhys 그대로의 수치 평가)
% ====================================================================
\begin{figure}[t]
\centering
\resizebox{\textwidth}{!}{%
\begin{tikzpicture}[>=Latex]
% =============== (a) 무차원 gap vs Omega/RT ===============
\begin{scope}[x=1.02cm,y=0.72cm]
% threshold band (single phase, gap identically 0)
\begin{scope}[on background layer]
  \fill[black!6] (0,-0.62) rectangle (2,5.62);
\end{scope}
\draw[->] (0,0) -- (6.55,0) node[right,font=\scriptsize] {$\Omega_j/RT$};
\draw[->] (0,-0.62) -- (0,5.62) node[above,font=\scriptsize] {$F\,\Delta U_j^\hys/RT$};
\foreach \xx in {1,2,3,4,5,6} {\draw (\xx,0.05) -- (\xx,-0.05) node[below,font=\scriptsize] {\xx};}
\foreach \yy in {1,2,3,4,5} {\draw (0.05,\yy) -- (-0.05,\yy) node[left,font=\scriptsize] {\yy};}
% threshold marker
\draw[densely dotted] (2,-0.62) -- (2,5.62);
\node[font=\scriptsize,align=center] at (1.0,3.1) {$\Omega_j\le2RT$:\\gap $\equiv0$\\(단상)};
% exact closed form (eq:dUhys), dimensionless
\draw[very thick] (0,0) -- (2,0);
\draw[very thick] plot[smooth] coordinates {(2,0) (2.02,0.0027) (2.05,0.0105) (2.1,0.0294)
  (2.2,0.0819) (2.35,0.1858) (2.5,0.3112) (2.75,0.5547) (3,0.8302) (3.25,1.13) (3.5,1.449)
  (3.75,1.7836) (4,2.1314) (4.5,2.8585) (5,3.6191) (5.5,4.4057) (6,5.2131)};
\node[font=\scriptsize,rotate=40] at (5.05,4.05) {정확식~\eqref{eq:dUhys}};
% soft cubic onset (8/3)u^3
\draw[densely dotted,thick] plot[smooth] coordinates {(2,0) (2.02,0.0026) (2.05,0.0102)
  (2.1,0.0277) (2.2,0.0731) (2.35,0.1533) (2.5,0.2385) (2.75,0.3798) (3,0.5132)
  (3.25,0.6361) (3.5,0.7482)};
\node[font=\scriptsize,anchor=west] at (3.28,0.42) {$\tfrac83u_j^3$ (연성 개시)};
% wrong constant-substitution Taylor: -(4/3)u^3
\draw[dashed,black!50] plot[smooth] coordinates {(2,0) (2.02,-0.0013) (2.05,-0.0051)
  (2.1,-0.0139) (2.2,-0.0365) (2.35,-0.0766) (2.5,-0.1193) (2.75,-0.1899) (3,-0.2566)
  (3.25,-0.318) (3.5,-0.3741)};
\node[font=\scriptsize,black!50,anchor=west] at (2.42,-0.95) {상수 대입 오답 $-\tfrac43u_j^3<0$};
\draw[->,thin,black!50] (3.05,-0.78) -- (3.28,-0.40);
% cross-check markers
\fill (3,0.8302) circle (1.1pt); \fill (4,2.1314) circle (1.1pt);
\node[font=\scriptsize,anchor=south east] at (3.06,0.90) {$(3,\,0.830)$};
\node[font=\scriptsize,anchor=south east] at (4.06,2.20) {$(4,\,2.131)$};
\node[font=\scriptsize] at (3.3,-1.55) {(a) 문턱$\cdot$연성 개시$\cdot$부호 함정};
\end{scope}
% =============== (b) gap(T) [mV], four staging Omega ===============
\begin{scope}[xshift=8.15cm,x=0.01235cm,y=0.0323cm,shift={(-240,0)}]
% operating window 260-340 K
\begin{scope}[on background layer]
  \fill[black!6] (260,0) rectangle (340,130);
\end{scope}
\draw[->] (240,0) -- (815,0) node[right,font=\scriptsize] {$T$ [K]};
\draw[->] (240,0) -- (240,130) node[above,font=\scriptsize] {$\Delta U_j^\hys(T)$ [mV]};
\foreach \xx in {300,400,500,600,700,800} {\draw (\xx,1.5) -- (\xx,-1.5) node[below,font=\scriptsize] {\xx};}
\foreach \yy in {25,50,75,100,125} {\draw (244,\yy) -- (236,\yy) node[left,font=\scriptsize] {\yy};}
\node[font=\tiny,black!55,anchor=north] at (321,127) {측정 창};
% 298.15 K guide
\draw[densely dotted] (298.15,0) -- (298.15,127);
% Omega = 13000 (2->1)
\draw[very thick] plot[smooth] coordinates {(240,125.375) (260,117.298) (280,109.546)
  (300,102.101) (320,94.95) (340,88.078) (360,81.477) (380,75.137) (400,69.049)
  (420,63.208) (440,57.609) (460,52.247) (480,47.118) (500,42.221) (520,37.553)
  (540,33.115) (560,28.906) (580,24.929) (600,21.185) (620,17.68) (640,14.42)
  (660,11.413) (680,8.671) (700,6.21) (720,4.055) (740,2.243) (760,0.84)
  (770,0.333) (777.8,0.066) (781.5,0.001)};
% Omega = 10000 (2L->2)
\draw[thick] plot[smooth] coordinates {(240,75.212) (260,68.272) (280,61.683)
  (300,55.432) (320,49.505) (340,43.894) (360,38.59) (380,33.589) (400,28.887)
  (420,24.483) (440,20.378) (460,16.576) (480,13.083) (500,9.911) (520,7.075)
  (540,4.6) (560,2.528) (580,0.932) (590,0.36) (597.4,0.075) (601.1,0.002)};
% Omega = 8000 (3->2L)
\draw[thick,black!55] plot[smooth] coordinates {(240,44.346) (260,38.459) (280,32.963)
  (300,27.849) (320,23.11) (340,18.743) (360,14.751) (380,11.142) (400,7.929)
  (420,5.137) (440,2.808) (460,1.024) (470,0.389) (477.1,0.084) (480.8,0.002)};
% Omega = 6000 (4->3)
\draw[thick,densely dashed] plot[smooth] coordinates {(240,17.332) (260,13.028)
  (280,9.226) (300,5.947) (320,3.23) (340,1.163) (350,0.433) (356.8,0.097) (360.5,0.002)};
% Tc ticks on axis (labels folded into curve tags below)
\foreach \tc in {360.8,481.1,601.4,781.8} {\draw[black!60,semithick] (\tc,0) -- (\tc,3.4);}
\draw[->,thin,black!60] (742,13.5) -- (778,1.5);
\node[font=\tiny,black!60,anchor=south] at (726,13.2) {눈금 $=T_{c,j}{=}\Omega_j/2R$};
% intersections at 298.15 K
\fill (298.15,102.78) circle (1.1pt) node[right,font=\tiny] {$102.8$};
\fill (298.15,56.00) circle (1.1pt) node[right,font=\tiny] {$56.0$};
\fill (298.15,28.31) circle (1.1pt) node[right,font=\tiny] {$28.3$};
\fill (298.15,6.23) circle (1.1pt) node[above right,font=\tiny] {$6.2$};
% curve labels (with their Tc)
\node[font=\tiny,anchor=west] at (556,32.5) {$\Omega{=}13000$ ($T_c\,782$ K)};
\node[font=\tiny,anchor=west] at (462,20.5) {$10000$ ($601$ K)};
\node[font=\tiny,anchor=west] at (390,14.2) {$8000$ ($481$ K)};
\node[font=\tiny,anchor=west] at (316,6.8) {$6000$ J/mol ($361$ K)};
% closing annotation
\node[font=\scriptsize,align=left,anchor=north west] at (480,122)
  {$T\!\to\!T_c{=}\Omega_j/2R$:\\ gap $\propto(T_c{-}T)^{3/2}$ 소멸\\ (도핑 $\Omega\!\downarrow$ 도 같은 축)};
\node[font=\scriptsize] at (520,-22) {(b) 온도$\cdot$$\Omega$ 에 따른 gap 소멸 (staging 초기값 4건)};
\end{scope}
\end{tikzpicture}%
}
\caption{히스테리시스 gap 닫힌 꼴 $\Delta U_j^\hys$(식~\eqref{eq:dUhys})의 전 거동 --- 좌표는 식
그대로의 수치 평가. (a) 무차원 gap $F\Delta U_j^\hys/RT=2[\omega u_j-2\,\mathrm{artanh}\,u_j]$
($\omega\equiv\Omega_j/RT$, $u_j=\sqrt{1-2/\omega}$): 문턱 $\Omega_j\le2RT$(음영)에서 정확히 $0$,
문턱 직후 두 급수를 함께 전개한 $\tfrac83u_j^3$(점선)으로 \emph{양수} 연성 개시하며, $\Omega_j$ 를
상수 대입만 한 잘못된 전개는 $-\tfrac43u_j^3$(회색 파선)로 부호가 뒤집힌다(\S\ref{sec:hys-gap} 의
Taylor 함정). 점 $(4,\,2.131)$ 은 그림~\ref{fig:hysloop} 의 $\Omega_j{=}4RT$ gap 과 같은 값이다.
(b) 표~\ref{tab:staging} 의 $\Omega_j$ 초기값 4건을 mV 로 평가한 $\Delta U_j^\hys(T)$: 각 곡선은
자기 임계온도 $T_{c,j}=\Omega_j/2R$(축 눈금)에서 $(T_c-T)^{3/2}$ 로 매끄럽게 소멸한다 --- 승온과
도핑에 의한 $\Omega_j$ 감소(식~\eqref{eq:lco-dope})가 같은 축 위에서 gap 을 닫는 두 손잡이다.
점선 수직선은 $298.15$ K, 그 교점 수치는 spinodal 상한 gap 이며 실측 분기는 $\gamma_j$ 배 안쪽이다
(식~\eqref{eq:Ubranch}).}
\label{fig:cand-ff1-1}
\end{figure}
